\section*{Capítulo \ref{ch:g11:minerals} --- Os minerais}

\begin{solution}{exo:g11:minerals:1}
A \omterm{def:g11:minerals:cleavage}{clivagem} é a tendência de um mineral a se partir ao longo de planos,
paralelos a direções em que as ligações entre os seus átomos são mais
fracas; a \omterm{def:g11:minerals:cleavage}{fratura} é a ruptura ao longo de superfícies irregulares, num
mineral sem plano de fraqueza. O quartzo não tem \omterm{def:g11:minerals:cleavage}{clivagem}: ele mostra uma
\omterm{def:g11:minerals:cleavage}{fratura} conchoidal, com superfícies lisas e curvas como as de um vidro
quebrado. Uma mica tem uma \omterm{def:g11:minerals:cleavage}{clivagem}, perfeita: ela se desfolha em lâminas.
\end{solution}

\begin{solution}{exo:g11:minerals:2}
Um traço de outro elemento basta para mudar a cor do quartzo, que pode ser
incolor, branco, preto, rosa, azul ou roxo; a olivina é quase sempre verde.
Mais confiáveis: a \omterm{def:g11:minerals:cleavage}{clivagem} (ou a \omterm{def:g11:minerals:cleavage}{fratura}), a dureza, o \omterm{def:g11:minerals:lustre}{brilho}, o \omterm{def:g11:minerals:habit}{hábito} dos
cristais livres, a densidade.
\end{solution}

\begin{solution}{exo:g11:minerals:3}
(1) Olivina. (2) Biotita. (3) Ortoclásio (um feldspato: duas clivagens,
dureza 6). (4) Quartzo (sem \omterm{def:g11:minerals:cleavage}{clivagem}, dureza 7).
\end{solution}

\begin{solution}{exo:g11:minerals:4}
(1) Cinco átomos: um silício no centro e quatro oxigênios nos vértices de um
tetraedro. (2) $+4 + 4 \times (-2) = -4$: o grupo é \ce{SiO4^4-}.
\end{solution}

\begin{solution}{exo:g11:minerals:5}
Olivina: tetraedros isolados. Piroxênios: cadeias simples. Anfibólios:
cadeias duplas. Micas: folhas. Feldspatos e quartzo: arcabouços.
\end{solution}

\begin{solution}{exo:g11:minerals:6}
(1) $M(\ce{SiO2}) = 28.085 + 2 \times 15.999 = \qty{60.083}{g/mol}$;
oxigênio: $31.998/60.083 \approx \qty{53.3}{\percent}$.
(2) $M(\ce{Mg2SiO4}) = 2 \times 24.305 + 28.085 + 4 \times 15.999 =
\qty{140.691}{g/mol}$; oxigênio: $63.996/140.691 \approx \qty{45.5}{\percent}$.
(3) A crosta, com \qty{46}{\percent}, é da mesma ordem: ela é feita quase
inteiramente de minerais ricos em oxigênio, os \omterm{def:g11:minerals:silicate}{silicatos}.
\end{solution}

\begin{solution}{exo:g11:minerals:7}
O grão branco tem duas clivagens e uma dureza entre 5.5 e 6.5: um feldspato
(6). O grão cinzento é mais duro que 6.5 e mostra uma \omterm{def:g11:minerals:cleavage}{fratura} curva,
conchoidal, sem \omterm{def:g11:minerals:cleavage}{clivagem}: quartzo (7).
\end{solution}

\begin{solution}{exo:g11:minerals:8}
(1) $4 - 2/2 = 3$ oxigênios por silício: \ce{SiO3}. (2) \ce{(Si,Al)2O6}
contém dois átomos de silício (ou de alumínio) para seis oxigênios, três por
silício: é \ce{SiO3} duas vezes.
\end{solution}

\begin{solution}{exo:g11:minerals:9}
$\rho = 0.55 \times 2.69 + 0.35 \times 3.4 + 0.10 \times 3.32 = 1.480 + 1.190
+ 0.332 \approx \qty{3.00}{g/cm^3}$. O cubo tem um volume de
\qty{1000}{cm^3}, logo uma massa de cerca de \qty{3.00}{kg}.
\end{solution}

\begin{solution}{exo:g11:minerals:10}
(1) $0.75 \times 3.32 + 0.25 \times 3.4 = 2.49 + 0.85 = \qty{3.34}{g/cm^3}$.
(2) $(3.34 - 3.00)/3.00 \approx 0.11$: cerca de \qty{11}{\percent} mais
denso. (3) O basalto e o \omterm{def:g11:ocean-ridges:gabbro}{gabro} da \omterm{def:g11:ocean-ridges:oceanic-lithosphere}{crosta oceânica} são mais leves que o
\omterm{def:g11:ocean-ridges:peridotite}{peridotito} do manto: eles ficam por cima dele.
\end{solution}

\begin{solution}{exo:g11:minerals:11}
(1) Plagioclásio $168/400 = 0.42$, piroxênio $0.33$, olivina $0.15$, buracos
$0.10$. (2) A proposição vale para uma rocha sem poros; aqui, um décimo do
volume está vazio. (3) Na parte sólida, as frações são $0.42/0.90
\approx 0.467$, $0.367$ e $0.167$:
$\rho_{\text{sólido}} \approx 0.467 \times 2.69 + 0.367 \times 3.4 + 0.167
\times 3.32 \approx \qty{3.06}{g/cm^3}$. Com os seus buracos, que
acrescentam volume e nenhuma massa, a rocha tem a densidade $0.42 \times 2.69 + 0.33 \times 3.4 +
0.15 \times 3.32 \approx \qty{2.75}{g/cm^3}$, isto é, $0.9 \times 3.06$.
\end{solution}

\begin{solution}{exo:g11:minerals:12}
(1) Oxigênio $464/16.0 = 29.0$~mol; silício $283/28.1 \approx 10.1$; alumínio
$84/27.0 \approx 3.1$; ferro $52/55.8 \approx 0.93$; cálcio $46/40.1 \approx
1.15$; magnésio $28/24.3 \approx 1.15$; sódio $23/23.0 = 1.0$; potássio
$15/39.1 \approx 0.38$. (2) Ao todo, cerca de $46.8$~mol de átomos, dos quais
$29.0/46.8 \approx 0.62$: quase dois átomos em cada três são de oxigênio.
(3) Os grandes íons oxigênio formam a maior parte dos átomos, e uma parte
ainda maior do volume: a crosta é, na escala dos átomos, um empilhamento de
íons oxigênio, com pequenos íons de silício e de metais nos espaços entre
eles.
\end{solution}

\begin{solution}{exo:g11:minerals:13}
(1) O número médio de vértices compartilhados é $(2 + 3)/2 = 2.5$, logo há
$4 - 2.5/2 = 2.75$ oxigênios por silício. (2) Quatro silícios para onze
oxigênios: \ce{Si4O11}. (3) \ce{Si7AlO22} contém oito tetraedros (sete de
silício, um de alumínio) para 22 oxigênios, isto é, $22/8 = 2.75$ por
tetraedro, ou \ce{Si4O11} duas vezes. Os dois grupos \ce{OH} não fazem
parte dos tetraedros.
\end{solution}

\begin{solution}{exo:g11:minerals:14}
$0.20 \times 2.62 + 0.10 \times 2.56 + 0.50 \times 2.65 + 0.10 \times 3.09 +
0.10 \times 3.23 = 0.524 + 0.256 + 1.325 + 0.309 + 0.323 \approx
\qty{2.74}{g/cm^3}$, acima de \qty{2.7}{g/cm^3}; o granito do exemplo, mais
rico em quartzo e em feldspato potássico, tem cerca de \qty{2.59}{g/cm^3},
abaixo de \qty{2.6}{g/cm^3}. As rochas do oeste são mais densas porque
contêm mais plagioclásio e mais biotita e hornblenda, ricas em ferro e
magnésio; os granitos ricos em sílica contêm mais quartzo e feldspato
potássico, os minerais formadores de rochas mais leves.
\end{solution}

\begin{solution}{exo:g11:minerals:15}
(1) Crosta superior: $66.6 \times 0.467 \approx \qty{31.1}{\percent}$ de
silício; crosta inteira: $60.6 \times 0.467 \approx \qty{28.3}{\percent}$.
(2) A crosta superior, mais rica em sílica, deve conter mais quartzo e
feldspatos, e menos minerais ricos em ferro e magnésio; a crosta mais
profunda, o contrário. (3) O quartzo e os feldspatos são os minerais
formadores de rochas mais leves: a crosta superior deve ser mais leve que a
crosta mais profunda.
\end{solution}

\begin{solution}{pb:g11:minerals:1}
\textbf{1.} Quartzo: sem \omterm{def:g11:minerals:cleavage}{clivagem}, \omterm{def:g11:minerals:cleavage}{fratura} conchoidal, \omterm{def:g11:minerals:lustre}{brilho} vítreo, mais
duro que o prego (7).

\textbf{2.} Ortoclásio, um feldspato potássico: duas clivagens, dureza 6 na
escala de Mohs.

\textbf{3.} Plagioclásio (aqui, oligoclásio), o feldspato com sódio e
cálcio.

\textbf{4.} D é a biotita, E a muscovita: duas micas. Os seus tetraedros
formam folhas, e as ligações entre as folhas são fracas: elas se partem ao
longo dessa única \omterm{def:g11:minerals:cleavage}{clivagem} perfeita.

\textbf{5.} O quartzo e os dois feldspatos (A, B, C) são arcabouços; a
biotita (D) contém ferro e magnésio.

\textbf{6.} Quartzo $150/500 = 0.30$; ortoclásio $0.35$; oligoclásio $0.25$;
biotita $0.06$; muscovita $0.04$.

\textbf{7.} Quando os grãos estão espalhados sem ordem, a parte dos pontos
sobre um mineral estima a parte da área que ele cobre na face, e esta é
igual à parte do volume que ele ocupa.

\textbf{8.} $0.35 + 0.25 = 0.60$: \qty{60}{\percent} de feldspato.

\textbf{9.} Em $1/500 = 0.002$, isto é, \qty{0.2}{\percent}.

\textbf{10.} Sim: feldspato, quartzo e mica.

\textbf{11.} $\rho = \phi_1 \rho_1 + \phi_2 \rho_2 + \dots + \phi_k \rho_k$,
com $\phi_1 + \dots + \phi_k = 1$.

\textbf{12.} Quartzo $0.30 \times 2.62 = 0.786$; ortoclásio $0.35 \times 2.56 =
0.896$; oligoclásio $0.25 \times 2.65 = 0.6625$; biotita $0.06 \times 3.09
\approx 0.185$; muscovita $0.04 \times 2.82 \approx 0.113$ (tudo em
\unit{g/cm^3}).

\textbf{13.} $0.786 + 0.896 + 0.6625 + 0.185 + 0.113 \approx
\qty{2.64}{g/cm^3}$.

\textbf{14.} O ortoclásio, com $0.896$, porque é o mais abundante. A biotita
e a muscovita, com \qty{3.09}{} e \qty{2.82}{g/cm^3}, são mais densas que os
outros minerais e elevam a densidade da rocha acima da do quartzo e do
feldspato sozinhos.

\textbf{15.} \qty{2640}{kg/m^3}: \qty{1}{m^3} do granito tem uma massa de
cerca de \qty{2.64}{t}.

\textbf{16.} $100 \times 2.643 \approx \qty{264.3}{g}$.

\textbf{17.} O bloco desloca \qty{100}{cm^3} de água, de massa
\qty{100}{g}: a balança indica $264.3 - 100 = \qty{164.3}{g}$.

\textbf{18.} $264.3/(264.3 - 164.3) = 264.3/100 \approx \qty{2.64}{g/cm^3}$.

\textbf{19.} Entre os dois: mais leve que a média de \qty{2.69}{g/cm^3},
acima dos \qty{2.6}{g/cm^3} dos granitos mais ricos em sílica e potássio.
Mais plagioclásio e mais biotita ou hornblenda, e menos quartzo e feldspato
potássico, tornariam um granito mais denso.

\textbf{20.} A densidade do granito é de cerca de
\textbf{\qty{2.64}{g/cm^3}}.
\end{solution}
